% **************************************************************************** %
\chapter{\code{LTspice} Schaltung eines Photovoltaik-Moduls}
\label{app:simu:module}
% **************************************************************************** %

\fref{fig:ltspice:solarCell}   zeigt   die  Implementation   unseres   Modells
einer     Solarzelle      gem\"ass     \fref{fig:circuit:solarCell}     (Seite
\pageref{fig:circuit:solarCell} in Abschnitt \ref{subsubsec:hw:ask:modell}) in
\code{LTspice}.

Ein MOSFET  wird benutzt, um  die Zelle  gesteuert (in unserer  Simulation bei
einer Tr\"agerfrequenz von \SI{10}{\kilo\hertz}) kurzschliessen zu k\"onnen.

Es    soll   das    f\"ur   den    MOSFET   schlimmere    Szenario   evaluiert
werden   (mehr    Strom/Leistung   durch    den   MOSFET\todo{Korrekt?}). Dazu
werden   aus  72   Zellen  zwei   verschiedene  Module   aufgebaut: Ein  Modul
\"ahnlich    zum    Sunset   Solargenerator    AS150    \cite{ref:solar:as150}
mit    ca. \SI{10}{\ampere}   Kurzschlussstrom    und   etwa    \SI{22}{\volt}
Leerlaufspannung  (\fref{fig:ltspice:module:cellBased:36x2}),  und  ein  Modul
\"ahnlich  zum Sunmodule  Pro-Series  XL SWB320  \cite{ref:solar:sunmodulePro}
mit   \SI{10}{\ampere}   Kurzschlussstrom    und   ungef\"ahr   \SI{45}{\volt}
Leerlaufspannung (\fref{fig:ltspice:module:cellBased:72x1}).

Weil die gesamte Kapazit\"at bei in Serie geschalteten Kondensatoren sinkt und
bei  paralleler  Anordnung steigt,  hat  die  $2 \times  36$-Parallelschaltung
\todo{Bindestrich und Spaces?} etwa die doppelte Kapazit\"at des Einzelstrangs
aus 72 Zellen. Der  durch den MOSFET fliessende Stroms, die  \"uber den MOSFET
abfallende Spannung und die im  MOSFET dissipierte Leistung beim Durchschalten
des Transistors sollen zwischen den beiden Szenarien verglichen werden.

Die zu  diesen Schaltungen geh\"orenden \code{.asc}-Dateien  sind elektronisch
verf\"ugbar  (Datentr\"ager  siehe Anhang  \ref{app:electronicStorage},  Seite
\pageref{app:electronicStorage}).

\begin{figure}[h!tb]
    \centering
    \includegraphics[width=\textwidth]{images/ltspice/singlecell.eps}
    \caption{%
        Schaltung        zur         Simulation        einer        Solarzelle
        gem\"ass        Abbildung       \ref{fig:circuit:solarCell}        von
        Seite     \pageref{fig:circuit:solarCell}. Die     Solarmodule     aus
        den     Abbildungen    \ref{fig:ltspice:module:cellBased:36x2}     und
        \ref{fig:ltspice:module:cellBased:72x1}  basieren auf  Zellen gem\"ass
        diesem Schaltschema.%
    }
    \label{fig:ltspice:solarCell}
\end{figure}

\begin{figure}[h!tb]
    \centering
    \includegraphics[width=\textwidth]{images/ltspice/module-72cells.eps}
    \caption{%
        Solarmodul     \"ahnlich     zu    Sunset     Solargenerator     AS150
        \cite{ref:solar:as150}, modelliert durch 2  parallele Str\"ange mit je
        36 Zellen gem\"ass Abbildung \ref{fig:ltspice:solarCell} in Serie.%
    }
    \label{fig:ltspice:module:cellBased:36x2}
\end{figure}

\begin{figure}[h!tb]
    \centering
    \includegraphics[width=\textwidth]{images/ltspice/module-72cells-series.eps}
    \caption{%
        Solarmodul   \"ahnlich   zu   Sunmodule  Pro-Series   XL   SW320   aus
        \cite{ref:solar:sunmodulePro},  modelliert durch  einen Strang  mit 72
        Zellen gem\"ass Abbildung \ref{fig:ltspice:solarCell} in Serie.%
    }
    \label{fig:ltspice:module:cellBased:72x1}
\end{figure}
\todo{Transistoren aus Schaltbildern entfernen, IN/OUT benennen, Abbildung sollte nur Modul sein}
